% !PRINTED 759
\weberpage{759}
%<<CONT
eine durch $p$ theilbare ganze rationale Zahl ist. Daraus ergiebt sich endlich, dass
\[
CF(0) + F(\alpha_1) + F(\alpha_2) + \cdots + F(\alpha_\mu)
\]
eine nicht durch $p$ theilbare, also auch nicht verschwindende, ganze rationale Zahl ist, die daher, vom Vorzeichen abgesehen, mindestens $=1$ ist.

Können wir nun endlich noch beweisen, dass auch bei der jetzigen Annahme über $f$ die Function $Q(x)$ für jedes endliche $x$ bei hinlänglicher Vergrösserung von $p$ beliebig klein gemacht werden kann, so ergiebt sich, genau wie oben, die Unmöglichkeit der Gleichung (4).

Dies lässt sich aber folgendermaassen einsehen: Wir betrachten statt der Function
\weq{6}{f(x) = c_nx^n + c_{n-1}x^{n-1} + \cdots + c_0}
die Function
\[
\begin{aligned}
\psi(x) &= \dfrac{a^{\mu p+p-1}x^{p-1}(x+a_1)^p(x+a_2)^p\ldots(x+a_\mu)^p}{\Pi(p-1)} \\
&= \gamma_nx^n + \gamma_{n-1}x^{n-1} + \cdots + \gamma_0,
\end{aligned}
\]
die nun lauter positive (aber nicht nothwendig rationale) Coëfficienten hat. Die Coëfficienten $c_n,c_{n-1}\ldots$ sind durch Multiplication und Addition aus den Zahlen $a,-\alpha_1,-\alpha_2,\ldots,-\alpha_\mu$ gebildet, und die entsprechenden Coëfficienten $\gamma_n,\gamma_{n-1}\ldots$ erhält man daraus, wenn man die $-\alpha_1,-\alpha_2,\ldots,-\alpha_\mu$ durch ihre absoluten Werthe $a_1,a_2,\ldots,a_\mu$ ersetzt, woraus nach dem schon oben erwähnten Satze der Einleitung folgt, dass die Coëfficienten $\gamma_n,\gamma_{n-1},\ldots$ gewiss nicht kleiner sind, als die absoluten Werthe der entsprechenden $c_n,c_{n-1}\ldots$

Nun ist für jedes endliche $x$, dessen absoluter Werth $\xi$ ist, der absolute Werth von $Q(x)$ nach \S.\,205, (10) kleiner, oder wenigstens nicht grösser als
\[
\gamma_n\xi^n + \gamma_{n-1}\xi^{n-1} + \cdots + \gamma_0 = \psi(\xi),
\]
und dass $\psi(\xi)$ unter jede Grenze heruntersinkt, wenn $p$ gross genug wird, schliesst man aus der Darstellung:
\[
\psi(x) = \dfrac{X}{ax}\dfrac{X^{p-1}}{\Pi(p-1)},
\]
wenn
\[
% [sic] -- the source prints Latin $u+1$; the derivation from the exponent $\mu p+p-1$ requires $\mu+1$.
X = a^{u+1}x(x+a_1)(x+a_2)\ldots(x+a_\mu)
\]
gesetzt ist.
